\section{Future work}
\label{sec:future_work}

This study moves the proof of concept onto \populace{}'s current full target surface, including
congressional-district targets. The next empirical step is not to hold districts back from the headline
fit, but to run separate robustness panels that ask narrower questions: which target families or
geographies generalize when withheld, which targets dominate the loss, and how sensitive the frontier is
to the production target weights and cap. Those panels should be read as diagnostics around the production
surface rather than as replacements for it.

A second thread is the build-large-then-prune regime the controls are designed for, in which a pipeline
assembles a candidate universe far larger than it can ship---millions of records rather than the 75{,}112
used here---and prunes it to a deployable budget under the calibration targets. The compression ratios
would be far more aggressive than the range tested here, which is where target-informed selection has its
clearest opportunity to justify its cost: a random draw from a tight budget is least likely to span a
demanding target system, so a selector trained on the targets has the most to add. Whether the current
$L_0$ implementation can realize that opportunity, and at what compute cost (Appendix~\ref{app:cost}), is
what that setting tests.

The full-surface design also clarifies which classical comparators are worth adding. Raking, GREG, balanced
sampling, and combinatorial optimization remain relevant reference methods, but not all are natural
baselines for a heterogeneous 31{,}534-target fiscal surface. The appropriate next comparison is on the
simpler categorical-margin subsets where their assumptions hold cleanly, in the spirit of
\citet{tanton2014co}'s comparison of generalized regression with combinatorial optimisation for
small-area reweighting, while the production frontier continues to use the Populace loss over the complete
surface.
